\SourceLeaf{386}%
sungen selbst verstatten, so dass man rückwärts aus jenen (den geodaetischen Längen und Breiten) die Winkelmessungen wieder wenigstens mit derselben Genauigkeit soll berechnen können, mit der sie angestellt sind, so sind jene abgekürzten Formeln bei weitem nicht zureichend, und bei sehr grossen Dreiecken muss man dann sogar wünschen, auch noch die Glieder der vierten Ordnung berücksichtigen zu können. Bei einer andern Form der Rechnung lässt sich diess durch sehr geschmeidige Methoden erreichen: es kann hier aber nicht der Ort sein, diess weiter zu entwickeln, und wir begnügen uns, diess Bedürfniss der höhern Geodaesie hier angedeutet zu haben. Das Werk selbst bietet verschiedene Fälle dar, wo die grossen Vortheile einer solchen Behandlungsweise fühlbar werden: so sind z.\,B. die auf der Turiner Sternwarte bei Gelegenheit der Azimuthalbestimmungen gemachten Einschneidungen der Dreieckspunkte Masse, Monte Soglio und Rocca Melone gar nicht benutzt, die unter jener Voraussetzung eine sehr schätzbare Controlle und Vergrösserung der Genauigkeit mit Leichtigkeit gegeben haben würden.

\Needspace{10\baselineskip}
\pdfbookmark[0]{Anzeigen. Mémorial du dépôt général de la guerre. 1830.}{anzeigen-memorial-1830}
\begin{center}
\rule{.85\textwidth}{.3pt}\par
Göttingische gelehrte Anzeigen.\quad 1830 April 29.\par
\rule{.85\textwidth}{.3pt}
\end{center}
\hspace*{\parindent}\textit{Mémorial du dépôt général de la guerre, imprimé par ordre du ministre.} T.\,I, 1829 (für 1802---1803), 696 S. T.\,III, 1826 (für 1825), 466 S. T.\,IV, 1828 (für 1826), 494 S. T.\,V, 1829 (für 1827---1828), 490 S. in 4. Nebst vielen Karten und Planen. Paris bei \textsc{Picquet}.

Die militärische Zeitschrift unter obigem Titel nahm im Jahre 1802, auf Veranlassung des Generals \textsc{Andreossy}, damaligen Directors des Depot, ihren Anfang. Man wollte nach und nach einen Theil der Schätze dieses grossartigen, während der Revolutionskriege ins unermessliche bereicherten Instituts, für die Kriegswissenschaft und die Hülfskenntnisse, Geschichte, Topographie, Geodaesie, Statistik u.\,s.\,w. gemeinnützig machen, und dazu die Musse des damals eingetretenen Friedens benutzen. So erschienen rasch nach einander die ersten Nummern dieser Zeitschrift; als jedoch der Krieg bald wieder ausbrach, gerieth der Fortgang derselben allmählich wieder ins Stocken, und mit der siebenten
